\section{The Stopping Recurrence}
\label{app:proofs}

This appendix states the dynamic program summarized in Section~\ref{sec:model}.
Nothing in the body depends on it. No theorem is claimed here: the recurrence is
standard backward induction, and its role in this paper is to be implemented exactly
and checked against enumeration, not to be new.

\subsection{Setup}

An episode has merchants $1, \dots, n$ in a fixed order, a hard budget vector $b$
over time, tokens, API calls, and spend, a hard purchase-price cap $P_{\max}$, and a
failure penalty $\Pi$ charged when the episode ends without a purchase. Querying
merchant $t$ consumes a known resource vector $c_t$ and reveals either that the offer
is unavailable, with weight $u_t$, or a price $p$ drawn from a finite support with
weights $w_{t,p}$. Let $W_t = u_t + \sum_p w_{t,p}$.

Offers are not recallable: an available offer is accepted when seen or lost.

\subsection{Recurrence}

Let $V_t(b)$ be the expected purchase loss before querying merchant $t$ with budget
$b$ remaining. Search ends when merchants run out or the next query is unaffordable:

\begin{equation}
  V_t(b) = \Pi
  \qquad\text{if } t > n \text{ or } c_t \not\le b .
\end{equation}

Otherwise the agent queries, observes, and either accepts or continues. Writing
$V_{t+1}(b - c_t)$ for the continuation value:

\begin{equation}
  V_t(b) = \frac{1}{W_t}\left[
    u_t\,V_{t+1}(b - c_t)
    + \sum_{p} w_{t,p}\,
      \begin{cases}
        \min\{p,\; V_{t+1}(b - c_t)\} & p \le P_{\max} \\
        V_{t+1}(b - c_t) & p > P_{\max}
      \end{cases}
  \right].
\end{equation}

The $\min$ is the decision: accept $p$ when it is no worse than continuing. An offer
above the hard cap cannot be bought at any price, so it contributes the continuation
value rather than $p$.

\subsection{The reservation price}

The optimal action is a threshold rule. Define

\begin{equation}
  R_t(b) = \min\{P_{\max},\; V_{t+1}(b - c_t)\},
\end{equation}

and accept the observed offer exactly when $p \le R_t(b)$. This is the reservation
price of Section~\ref{sec:model}: the cap, or the value of continuing, whichever
binds first.

The three qualitative forces described in the body follow directly. When $t = n$ or
no further query is affordable, $V_{t+1} = \Pi$ and the threshold rises to the cap,
so the agent accepts. Raising $u_t$ mixes more of the fallback into the continuation
value and raises the threshold. Making a future merchant cheaper lowers $V_{t+1}$ and
lowers the threshold.

\subsection{Deriving the closed form}

Section~\ref{sec:rule} states the rule without derivation. Take prices i.i.d.\
uniform on $[a, b]$, ignore stockouts, and let every query cost the same, so the
budget fixes a horizon $k$ through \eqref{eq:horizon}. Then $R_k$ depends only on
$k$, and the threshold identity is

\begin{equation}
  R_k = \mathbb{E}\!\left[\min(P, R_{k-1})\right],
  \qquad R_0 = b .
\end{equation}

For $r \in [a,b]$ and $P \sim \mathrm{Unif}[a,b]$,

\begin{equation}
  \mathbb{E}[\min(P, r)]
  = r - \frac{(r-a)^2}{2(b-a)} .
\end{equation}

Substituting $u_k = (R_k - a)/(b-a)$ and cancelling $(b-a)$ leaves the
one-dimensional map

\begin{equation}
  u_k = u_{k-1} - \tfrac{1}{2} u_{k-1}^2, \qquad u_0 = 1,
\end{equation}

which is \eqref{eq:closedform}. The first terms are $u_1 = \tfrac12$,
$u_2 = \tfrac38$, $u_3 = \tfrac{39}{128}$.

Two remarks. First, the map is exactly the Cayley--Moser recursion: writing
$\mu_k = 1 - u_k$ gives $\mu_k = \tfrac12(1 + \mu_{k-1}^2)$, the standard
full-information form~\citep{gilbert1966recognizing,ferguson1989secretary}. The
asymptotic $u_k \sim 2/k$ follows from the logistic form
$\alpha_k = \alpha_{k-1}(1 - \alpha_{k-1})$ under $u_k = 2\alpha_k$, for which a
full expansion is known~\citep{finch2024quadratic}. Second, no inspection cost
appears anywhere above; cost is absent from the objective and enters only through
the horizon.

Since $u_k$ is decreasing in $k$, the threshold rises as the budget drains, and at
$k = 0$ it reaches $b$: with nothing reachable left the agent accepts any admissible
offer rather than incur the failure penalty.

We verify the sequence numerically against the solver rather than relying on the
algebra alone; agreement is within \ResultClosedFormMaxRelativeErrorPercent\% across
the horizons checked, the residual being the discrete price grid the solver requires.

\subsection{Adaptive routing}

When the agent may also choose \emph{which} merchant to query, the state carries the
set $S$ of remaining merchants and the recurrence minimizes over the choice:

\begin{equation}
  V(S, b) = \min_{t \in S,\; c_t \le b} \; Q_t(S, b),
\end{equation}

where $Q_t(S, b)$ is the bracketed expectation above with $V_{t+1}(b - c_t)$ replaced
by $V(S \setminus \{t\}, b - c_t)$, and $V(S,b) = \Pi$ when no $t \in S$ is
affordable. Ties are broken by lowest index so the policy is deterministic.

\subsection{Computation}

All arithmetic is exact and rational; no floating-point tolerance is involved.
Subproblems are memoized on $(S, b)$ or $(t, b)$. The fixed-order program is linear
in $n$ for a given budget path, while the routing version enumerates subsets and is
exponential in $n$, which is acceptable at the horizons in this paper and is a real
obstacle beyond them.

\subsection{Verification}

For instances small enough to enumerate every deterministic accept-or-continue
policy, the minimum enumerated expected loss is compared against $V_1(b)$ in exact
arithmetic. The registered instances cover distinct prices, stockout weight, a
binding price cap, and a penalty large enough to dominate the price term. All
\ResultSolverCasesVerified{} agree exactly (claim \texttt{SOLVER-AGREEMENT-001}).
Agreement on these instances does not prove optimality in general, and the
fixed-order instances say nothing about the routing variant.
